\chapter{सीमाएँ और सांतत्य}\label{ch:g12:limcont}

यह अध्याय सीमा की धारणा को \omterm{def:g11:seq:sequence}{अनुक्रमों} से बढ़ाकर वास्तविक चर के \omterm{def:g10:functions:function}{फलनों} तक ले
जाता है, और \omterm{def:g12:limcont:continuity}{सांतत्य} प्रस्तुत करता है। केंद्रीय परिणाम मध्यवर्ती मान प्रमेय
है, जो गुणात्मक सूचना (``$f$ \omterm{def:g12:limcont:continuity}{संतत} है और चिह्न बदलता है'') को \omterm{def:g10:algebra:equation}{समीकरणों} के हलों
के अस्तित्व में बदल देती है।

\section{फलनों की सीमाएँ}

पूरे अध्याय में $f$ किसी \omterm{def:g10:numbers:interval}{अंतराल} या \omterm{def:g10:numbers:interval}{अंतरालों} के \omterm{def:g10:numbers:interunion}{सम्मिलन} $D \subseteq \R$ पर परिभाषित \omterm{def:g11:func:function}{फलन}
है।

\begin{definition}[अनंत पर सीमा]\label{def:g12:limcont:liminf}
मान लीजिए $f$ किसी \omterm{def:g10:numbers:interval}{अंतराल} $\intco{a}{+\infty}$ पर परिभाषित है और $\ell \in \R$ है। हम कहते हैं कि $f$
\emph{$+\infty$ पर $\ell$ की ओर जाता है}\index{सीमा!फलन की}, और $\lim\limits_{x\to+\infty} f(x) = \ell$ लिखते हैं, यदि
$\ell$ को समेटने वाले हर खुले \omterm{def:g10:numbers:interval}{अंतराल} में पर्याप्त बड़े $x$ के लिए सभी मान $f(x)$
आ जाएँ: अर्थात् हर $\varepsilon > 0$ के लिए ऐसा $A$ हो कि सभी $x \geq A$ के लिए $\abs{f(x) - \ell} \leq \varepsilon$ हो।

हम कहते हैं कि $f$ $+\infty$ पर $+\infty$ की ओर जाता है यदि हर $M \in \R$ के लिए ऐसा $A$
हो कि सभी $x \geq A$ के लिए $f(x) \geq M$ हो। $-\infty$ पर सीमाएँ और $-\infty$ के बराबर सीमाएँ इसी
तरह परिभाषित होती हैं।
\end{definition}

\begin{definition}[किसी बिंदु पर सीमा]\label{def:g12:limcont:limpoint}
मान लीजिए $a \in \R$ है और $f$ $a$ के पास (संभवतः $a$ को छोड़कर) परिभाषित है।
हम कहते हैं कि $f$ $a$ पर $\ell \in \R$ की ओर जाता है, और $\lim\limits_{x \to a} f(x) = \ell$ लिखते हैं, यदि हर
$\varepsilon > 0$ के लिए ऐसा $\delta > 0$ हो कि $\abs{x - a} \leq \delta$ वाले सभी $x \in D$ के लिए $\abs{f(x) - \ell} \leq \varepsilon$ हो।

हम कहते हैं कि $f$ $a$ पर $+\infty$ की ओर जाता है यदि हर $M$ के लिए ऐसा $\delta > 0$ हो
कि जब भी $x \in D$ और $\abs{x - a} \leq \delta$ हों तब $f(x) \geq M$ हो। एकपक्षीय सीमाएँ ($x \to a^+$, $x \to a^-$) $x$ को
$x > a$ या $x < a$ तक सीमित कर देती हैं।
\end{definition}

\begin{example}
$\lim\limits_{x\to+\infty} \frac{1}{x} = 0$, $\lim\limits_{x\to 0^+} \frac{1}{x} = +\infty$, $\lim\limits_{x\to 0^-} \frac{1}{x} = -\infty$। \omterm{def:g11:func:function}{फलन} $x \mapsto \frac1x$ की $0$ पर कोई (द्विपक्षीय) सीमा नहीं है।
\end{example}

\begin{definition}[अनन्तस्पर्शी]\label{def:g12:limcont:asymptote}
रेखा $y = \ell$ $+\infty$ पर (क्रमशः $-\infty$ पर) $f$ के वक्र की \emph{क्षैतिज
अनन्तस्पर्शी}\index{अनन्तस्पर्शी} है यदि $\lim_{x\to+\infty} f(x) = \ell$ हो (क्रमशः $-\infty$ पर)। रेखा $x = a$
\emph{ऊर्ध्वाधर अनन्तस्पर्शी}\index{अनन्तस्पर्शी} है यदि $a$ पर $f$ की कोई
एकपक्षीय सीमा अनंत हो।
\end{definition}

\begin{omfigure}
\begin{tikzpicture}
\begin{axis}[omaxis,
  width=8.8cm, height=6cm,
  xmin=-2.9, xmax=4.9, ymin=-5.2, ymax=9.2,
  xtick={1}, ytick={2}]
  \draw[omProp, dashed] (axis cs:-2.9,2) -- (axis cs:4.9,2);
  \draw[omProp, dashed] (axis cs:1,-5.2) -- (axis cs:1,9.2);
  \addplot[omDef, thick, domain=-2.8:0.855, samples=60] {2 + 1/(x-1)};
  \addplot[omDef, thick, domain=1.145:4.8, samples=60] {2 + 1/(x-1)};
\end{axis}
\end{tikzpicture}

{\small $f(x) = 2 + \frac{1}{x-1}$ का वक्र, उसकी \omterm{def:g12:limcont:asymptote}{क्षैतिज अनन्तस्पर्शी} $y = 2$ ($\pm\infty$ पर) और \omterm{def:g12:limcont:asymptote}{ऊर्ध्वाधर
अनन्तस्पर्शी} $x = 1$ के साथ।}
\end{omfigure}

\begin{proposition}[सीमाओं पर संक्रियाएँ]\label{prop:g12:limcont:operations}
\cref{prop:g12:seq:operations} के नियम (योग, गुणनफल, भागफल) \omterm{def:g10:functions:function}{फलनों} की सीमाओं के लिए भी ज्यों के त्यों
लागू होते हैं, चाहे किसी बिंदु पर हों या अनंत पर, और वही चार अनिर्धार्य रूप
$(+\infty)+(-\infty)$, $0\times\infty$, $\frac{\infty}{\infty}$, $\frac{0}{0}$ रहते हैं।
\end{proposition}

\begin{proof}
उपपत्तियाँ \omterm{def:g11:seq:sequence}{अनुक्रमों} वाली स्थिति जैसी ही हैं, केवल ``$n \geq N$ के लिए'' के स्थान
पर ``$x \geq A$ के लिए'' या ``$\abs{x-a} \leq \delta$ के लिए'' रख दीजिए।
\end{proof}

\begin{theorem}[संयोजन की सीमा]\label{thm:g12:limcont:composition}
मान लीजिए $f, g$ \omterm{def:g11:func:function}{फलन} हैं और $a, b, c$ \omterm{def:g10:numbers:sets}{वास्तविक संख्याएँ} या $\pm\infty$ हैं। यदि $\lim\limits_{x \to a} f(x) = b$ और
$\lim\limits_{y \to b} g(y) = c$ हों, तो
\[
\lim_{x \to a} g\bigl(f(x)\bigr) = c .
\]
यही कथन $f$ के स्थान पर किसी \omterm{def:g12:seq:sequence}{अनुक्रम} के साथ भी सही है: यदि $u_n \to b$ और $\lim_{y\to b} g(y) = c$ हों,
तो $g(u_n) \to c$।
\end{theorem}

\begin{proof}
$b, c$ को परिमित लीजिए (शेष स्थितियाँ इसी जैसी हैं)। मान लीजिए $\varepsilon > 0$ है। ऐसा
$\delta > 0$ है कि $\abs{y - b} \leq \delta$ से $\abs{g(y) - c} \leq \varepsilon$ निकलता है; और ऐसा $\delta' > 0$ है कि $\abs{x - a} \leq \delta'$ से $\abs{f(x) - b} \leq \delta$ निकलता
है। दोनों को जोड़ने पर $\abs{x - a} \leq \delta'$ के लिए $\abs{g(f(x)) - c} \leq \varepsilon$ मिलता है।
\end{proof}

\begin{theorem}[तुलना और संपीडन]\label{thm:g12:limcont:squeeze}
मान लीजिए $f, g, h$ \omterm{def:g11:func:function}{फलन} हैं और $a \in \R \cup \{\pm\infty\}$ है।
\begin{enumerate}
  \item यदि $a$ के पास $f \leq g$ हो और $x \to a$ पर $f(x) \to +\infty$ हो, तो $g(x) \to +\infty$।
  \item यदि $a$ के पास $f \leq g \leq h$ हो और $f, h$ $a$ पर एक ही परिमित सीमा $\ell$
        की ओर जाएँ, तो $x \to a$ पर $g(x) \to \ell$।
\end{enumerate}
\end{theorem}

\begin{proof}
वही तर्क जो \omterm{def:g11:seq:sequence}{अनुक्रमों} के लिए था (\cref{thm:g12:seq:squeeze})।
\end{proof}

\begin{method}[व्यवहार में सीमाएँ निकालना]\label{met:g12:limcont:limits}
\begin{itemize}
  \item $\pm\infty$ पर बहुपद और परिमेय \omterm{def:g11:func:function}{फलन}: $x$ की सबसे बड़ी घात बाहर निकाल
        लीजिए; सीमा अग्र पदों के अनुपात की सीमा होती है।
  \item किसी बिंदु $a$ पर $\frac{0}{0}$: अंश और हर दोनों से $(x-a)$ बाहर निकालिए,
        या किसी \omterm{def:g11:deriv:derivative}{अवकलज} $\lim_{x\to a}\frac{f(x)-f(a)}{x-a} = f'(a)$ को पहचानिए (\cref{ch:g12:deriv})।
  \item वर्गमूल: संयुग्मी व्यंजक से गुणा कीजिए।
  \item दोलन करते गुणनखंड ($\cos$, $\sin$): संपीडन प्रमेय काम में लाइए।
\end{itemize}
\end{method}

\section{सांतत्य}

\begin{definition}[सांतत्य]\label{def:g12:limcont:continuity}
मान लीजिए $f$ किसी \omterm{def:g10:numbers:interval}{अंतराल} $I$ पर परिभाषित है और $a \in I$ है। \omterm{def:g11:func:function}{फलन} $f$
\emph{$a$ पर संतत}\index{सांतत्य} है यदि $\lim\limits_{x \to a} f(x) = f(a)$ हो। वह \emph{$I$ पर
संतत}\index{संतत} है यदि वह $I$ के हर बिंदु पर संतत हो।
\end{definition}

अंतर्ज्ञान से: किसी \omterm{def:g10:numbers:interval}{अंतराल} पर \omterm{def:g12:limcont:continuity}{संतत} \omterm{def:g11:func:function}{फलन} का \omterm{def:g10:functions:graph}{आलेख} क़लम उठाए बिना खींचा जा सकता
है।

\begin{proposition}[सामान्य फलनों का सांतत्य]\label{prop:g12:limcont:usual}
बहुपद, परिमेय \omterm{def:g11:func:function}{फलन}, $x \mapsto \sqrt{x}$, $x \mapsto \abs{x}$, $\cos$, $\sin$, $\exp$ और $\ln$ अपने \omterm{def:g11:func:function}{प्रांत} में
समाए हर \omterm{def:g10:numbers:interval}{अंतराल} पर \omterm{def:g12:limcont:continuity}{संतत} हैं। \omterm{def:g12:limcont:continuity}{संतत} \omterm{def:g10:functions:function}{फलनों} के योग, गुणनफल, भागफल (जहाँ परिभाषित
हों) और संयोजन भी \omterm{def:g12:limcont:continuity}{संतत} होते हैं।
\end{proposition}

\begin{proof}[आंशिक उपपत्ति]
योग, गुणनफल, भागफल और संयोजन के अंतर्गत बंदपन सीमाओं की संगत प्रमेयों से
निकलता है। $x \mapsto x$ और अचर \omterm{def:g10:functions:function}{फलनों} का \omterm{def:g12:limcont:continuity}{सांतत्य} परिभाषा से तत्काल मिल जाता है, इसलिए
बहुपद और परिमेय \omterm{def:g11:func:function}{फलन} \omterm{def:g12:limcont:continuity}{संतत} हैं। शेष सामान्य \omterm{def:g10:functions:function}{फलनों} का \omterm{def:g12:limcont:continuity}{सांतत्य} उन्हीं अध्यायों
में सिद्ध होता है जहाँ उनकी रचना होती है, या स्वीकार कर लिया जाता है।
\end{proof}

\begin{example}
\emph{महत्तम \omterm{def:g10:numbers:sets}{पूर्णांक} \omterm{def:g11:func:function}{फलन}} $x \mapsto \floor{x}$ ($\leq x$ वाला सबसे बड़ा \omterm{def:g10:numbers:sets}{पूर्णांक}) किसी भी
\omterm{def:g10:numbers:sets}{पूर्णांक} $n$ पर \omterm{def:g12:limcont:continuity}{संतत} नहीं है: वहाँ उसकी बाईं सीमा $n - 1$ है जबकि $\floor{n} = n$।
\end{example}

\begin{omfigure}
\begin{tikzpicture}
\begin{axis}[omaxis,
  width=7.2cm, height=5.6cm,
  xmin=-1.9, xmax=3.5, ymin=-2.5, ymax=2.7,
  xtick={-1,1,2,3}, ytick={-2,-1,1,2}]
  \draw[omDef, thick] (axis cs:-1.5,-2) -- (axis cs:-1,-2);
  \draw[omDef, thick] (axis cs:-1,-1) -- (axis cs:0,-1);
  \draw[omDef, thick] (axis cs:0,0) -- (axis cs:1,0);
  \draw[omDef, thick] (axis cs:1,1) -- (axis cs:2,1);
  \draw[omDef, thick] (axis cs:2,2) -- (axis cs:3,2);
  \fill[omDef] (axis cs:-1,-1) circle (1.5pt);
  \fill[omDef] (axis cs:0,0) circle (1.5pt);
  \fill[omDef] (axis cs:1,1) circle (1.5pt);
  \fill[omDef] (axis cs:2,2) circle (1.5pt);
  \draw[omDef, fill=white] (axis cs:-1,-2) circle (1.5pt);
  \draw[omDef, fill=white] (axis cs:0,-1) circle (1.5pt);
  \draw[omDef, fill=white] (axis cs:1,0) circle (1.5pt);
  \draw[omDef, fill=white] (axis cs:2,1) circle (1.5pt);
  \draw[omDef, fill=white] (axis cs:3,2) circle (1.5pt);
\end{axis}
\end{tikzpicture}

{\small महत्तम \omterm{def:g10:numbers:sets}{पूर्णांक} \omterm{def:g11:func:function}{फलन} हर \omterm{def:g10:numbers:sets}{पूर्णांक} पर उछल जाता है: मान (भरा बिंदु) बाईं
सीमा (खोखला बिंदु) से भिन्न है।}
\end{omfigure}

\begin{proposition}[अनुक्रम और सांतत्य]\label{prop:g12:limcont:seqcont}
यदि $u_n \to \ell$ हो और $f$ $\ell$ पर \omterm{def:g12:limcont:continuity}{संतत} हो, तो $f(u_n) \to f(\ell)$। विशेष रूप से, यदि $u_{n+1} = f(u_n)$ से
परिभाषित कोई \omterm{def:g12:seq:sequence}{अनुक्रम} $\ell$ पर \omterm{def:g12:seq:limit}{अभिसरित} हो और $f$ $\ell$ पर \omterm{def:g12:limcont:continuity}{संतत} हो, तो $f(\ell) = \ell$।
\end{proposition}

\begin{proof}
पहला दावा \cref{thm:g12:limcont:composition} को \omterm{def:g12:seq:sequence}{अनुक्रम} $(u_n)$ के साथ लगाने पर मिलता है। दूसरे के लिए $u_{n+1} = f(u_n)$
के दोनों पक्षों में सीमा लीजिए: बायाँ पक्ष $\ell$ की ओर जाता है, दाहिना $f(\ell)$
की ओर, और सीमाएँ अद्वितीय होती हैं।
\end{proof}

\section{मध्यवर्ती मान प्रमेय}

\begin{theorem}[मध्यवर्ती मान प्रमेय]\label{thm:g12:limcont:ivt}
\index{मध्यवर्ती मान प्रमेय}
मान लीजिए $f$ $\intcc{a}{b}$ पर \omterm{def:g12:limcont:continuity}{संतत} है और $k$ $f(a)$ तथा $f(b)$ के बीच की कोई
संख्या है। तब ऐसा $c \in \intcc{a}{b}$ है कि $f(c) = k$ हो।
\end{theorem}

\begin{omfigure}
\begin{tikzpicture}
\begin{axis}[omaxis,
  width=8.8cm, height=6cm,
  xmin=-0.4, xmax=4.4, ymin=-0.5, ymax=4.7,
  xtick={0.5,3.19,3.8}, xticklabels={$a$,$c$,$b$},
  ytick={1.26,2.5,3.98}, yticklabels={$f(a)$,$k$,$f(b)$}]
  \addplot[omProp, dashed, domain=0:3.19] {2.5};
  \addplot[omDef, thick, domain=0.5:3.8] {0.08*x^3 - 0.5*x + 1.5};
  \draw[black, dashed, thin] (axis cs:0.5,0) -- (axis cs:0.5,1.26)
    -- (axis cs:0,1.26);
  \draw[black, dashed, thin] (axis cs:3.8,0) -- (axis cs:3.8,3.98)
    -- (axis cs:0,3.98);
  \draw[black, dashed, thin] (axis cs:3.19,2.5) -- (axis cs:3.19,0);
  \fill (axis cs:0.5,1.26) circle (1.5pt);
  \fill (axis cs:3.8,3.98) circle (1.5pt);
  \fill (axis cs:3.19,2.5) circle (1.5pt);
\end{axis}
\end{tikzpicture}

{\small मध्यवर्ती मान प्रमेय: $(a, f(a))$ को $(b, f(b))$ से जोड़ने वाले किसी भी \omterm{def:g12:limcont:continuity}{संतत} वक्र
को $f(a)$ और $f(b)$ के बीच पड़ी हर क्षैतिज रेखा $y = k$ को पार करना ही पड़ता है।}
\end{omfigure}

\begin{proof}
$f$ के स्थान पर $f - k$ रखकर (और आवश्यकता हो तो $k - f$ से बदलकर) हम मान सकते
हैं कि $f(a) \leq 0 \leq f(b)$ है और $f$ का शून्य ढूँढ़ रहे हैं। हम
\emph{द्विभाजन}\index{द्विभाजन} से दो \omterm{def:g12:seq:sequence}{अनुक्रम} बनाते हैं: $a_0 = a$, $b_0 = b$ रखिए;
$f(a_n) \leq 0 \leq f(b_n)$ वाले $a_n \leq b_n$ दिए होने पर $m = \frac{a_n + b_n}{2}$ लीजिए और परिभाषित कीजिए
\[
(a_{n+1}, b_{n+1}) =
\begin{cases}
(a_n, m) & \text{यदि } f(m) \geq 0,\\
(m, b_n) & \text{यदि } f(m) < 0.
\end{cases}
\]
दोनों स्थितियों में $f(a_{n+1}) \leq 0 \leq f(b_{n+1})$ रहता है, \omterm{def:g12:seq:sequence}{अनुक्रम} $(a_n)$ \omterm{def:g12:seq:monotonic}{वर्धमान} है, $(b_n)$ \omterm{def:g10:functions:variations}{ह्रासमान}
है, और $b_n - a_n = \frac{b-a}{2^n} \to 0$: \omterm{def:g12:seq:sequence}{अनुक्रम} संलग्न हैं (\cref{exo:g12:seq:10} देखिए) और एक उभयनिष्ठ सीमा $c \in \intcc{a}{b}$ पर
\omterm{def:g12:seq:limit}{अभिसरित} होते हैं। \omterm{def:g12:limcont:continuity}{सांतत्य} से $f(a_n) \to f(c)$ और $f(b_n) \to f(c)$; $f(a_n) \leq 0$ तथा $f(b_n) \geq 0$ में सीमा लेने
पर $f(c) \leq 0 \leq f(c)$ मिलता है, इसलिए $f(c) = 0$।
\end{proof}

\begin{remark}
\omterm{thm:g12:limcont:ivt}{द्विभाजन} वाली उपपत्ति एक कलनविधि है: वह किसी हल के मनचाहे यथार्थ \omterm{def:g10:numbers:approx}{सन्निकटन} दे
देती है और हर क़दम पर त्रुटि आधी कर देती है। संगणक केवल $f$ का मान निकालना
जानते हुए भी $f(x) = 0$ इसी तरह हल कर लेता है।
\end{remark}

\begin{theorem}[एकैकी आच्छादन प्रमेय]\label{thm:g12:limcont:bijection}
मान लीजिए $f$ किसी \omterm{def:g10:numbers:interval}{अंतराल} $I$ पर \omterm{def:g12:limcont:continuity}{संतत} और \emph{निरंतर \omterm{def:g12:seq:monotonic}{एकदिष्ट}} है। तब
$I$ के सिरों पर $f$ की सीमाओं के कड़ाई से बीच पड़ने वाले हर $k$ के लिए
\omterm{def:g10:algebra:equation}{समीकरण} $f(x) = k$ का $I$ में \emph{अद्वितीय} हल है।
\end{theorem}

\begin{proof}
अस्तित्व: ऐसे $a, b \in I$ चुनिए कि $k$ $f(a)$ और $f(b)$ के बीच हो (यह संभव है
क्योंकि $f$ के मान उसकी सिरा-सीमाओं के पास पहुँचते हैं), और $\intcc{a}{b}$ पर
मध्यवर्ती मान प्रमेय लगाइए। अद्वितीयता: यदि $x < y$ हो, तो निरंतर एकदिष्टता से
$f(x) \neq f(y)$ मिलता है, इसलिए दो भिन्न बिंदु दोनों हल नहीं हो सकते।
\end{proof}

\begin{method}[एकैकी आच्छादन प्रमेय से $f(x) = k$ हल करना]\label{met:g12:limcont:bijection}
यह सिद्ध करने के लिए कि किसी \omterm{def:g10:numbers:interval}{अंतराल} में $f(x) = k$ का ठीक एक हल है:
\begin{enumerate}
  \item $f$ की \omterm{def:g10:functions:table}{परिवर्तन-सारणी} बनाइए ($f'$ का चिह्न, \cref{ch:g12:deriv} देखिए);
  \item हर उस \omterm{def:g10:numbers:interval}{अंतराल} पर जहाँ $f$ \omterm{def:g12:limcont:continuity}{संतत} और निरंतर \omterm{def:g12:seq:monotonic}{एकदिष्ट} है, $k$ की तुलना
        सिरों के मानों या सीमाओं से कीजिए;
  \item ऐसे हर \omterm{def:g10:numbers:interval}{अंतराल} पर एकैकी आच्छादन प्रमेय लगाइए और हल गिन लीजिए।
\end{enumerate}
इसके बाद कैलकुलेटर या \omterm{thm:g12:limcont:ivt}{द्विभाजन} हर हल को जितना चाहें उतना बारीकी से बता देता
है।
\end{method}

\begin{example}\label{ex:g12:limcont:cubic}
\omterm{def:g10:algebra:equation}{समीकरण} $x^3 + x = 1$ का ठीक एक वास्तविक हल है। सचमुच $f(x) = x^3 + x$ $\R$ पर \omterm{def:g12:limcont:continuity}{संतत} और निरंतर
\omterm{def:g12:seq:monotonic}{वर्धमान} है (निरंतर \omterm{def:g11:func:monotone}{वर्धमान फलनों} के योग के रूप में), $\lim_{-\infty} f = -\infty$ और $\lim_{+\infty} f = +\infty$, इसलिए
$k = 1$ के साथ एकैकी आच्छादन प्रमेय लागू होती है। चूँकि $f(0.6) = 0.816 < 1$ और $f(0.7) = 1.043 > 1$ हैं,
इसलिए हल $\intoo{0.6}{0.7}$ में है।
\end{example}

\section{अभ्यास}

\begin{exercise}[$\star$]\label{exo:g12:limcont:1}
निम्नलिखित सीमाएँ निकालिए:
\[
\lim_{x\to+\infty} \frac{2x^3 - x + 1}{5x^3 + x^2}, \qquad
\lim_{x\to-\infty} \bigl(x^2 + 3x - 1\bigr), \qquad
\lim_{x\to 2^+} \frac{x+1}{x-2}, \qquad
\lim_{x\to+\infty} \frac{\cos x}{x}.
\]
\end{exercise}

\begin{exercise}[$\star$]\label{exo:g12:limcont:2}
मान लीजिए $x \neq 1$ के लिए परिभाषित $f(x) = \dfrac{2x^2 - 3x + 1}{x - 1}$ है।
\begin{enumerate}
  \item दिखाइए कि सभी $x \neq 1$ के लिए $f(x) = 2x - 1$ है।
  \item $\lim\limits_{x \to 1} f(x)$ निष्कर्ष निकालिए। क्या $f$ को $1$ पर \omterm{def:g12:limcont:continuity}{संतत} \omterm{def:g11:func:function}{फलन} तक बढ़ाया
        जा सकता है?
\end{enumerate}
\end{exercise}

\begin{exercise}[$\star$]\label{exo:g12:limcont:3}
$f(x) = \dfrac{3x - 1}{x + 2}$ के वक्र की अनन्तस्पर्शियाँ ज्ञात कीजिए और वक्र का रेखाचित्र बनाइए।
\end{exercise}

\begin{exercise}[$\star\star$]\label{exo:g12:limcont:4}
निकालिए:
\[
\lim_{x\to+\infty} \bigl(\sqrt{x^2 + x} - x\bigr)
\qquad\text{और}\qquad
\lim_{x\to 0} \frac{\sqrt{1+x} - 1}{x}.
\]
\end{exercise}

\begin{exercise}[$\star\star$]\label{exo:g12:limcont:5}
मान लीजिए $f(x) = x + \sin x$ है।
\begin{enumerate}
  \item दिखाइए कि $x \to +\infty$ पर $f(x) \to +\infty$ है, यद्यपि $f$ किसी भी \omterm{def:g10:numbers:interval}{अंतराल} $\intco{A}{+\infty}$ पर
        \omterm{def:g12:seq:monotonic}{एकदिष्ट} नहीं है।
  \item दिखाइए कि $g(x) = \dfrac{x + \sin x}{x - \sin x}$ $+\infty$ पर $1$ की ओर जाता है।
\end{enumerate}
\end{exercise}

\begin{exercise}[$\star\star$]\label{exo:g12:limcont:6}
दिखाइए कि \omterm{def:g10:algebra:equation}{समीकरण} $x^3 - 3x + 1 = 0$ के ठीक तीन वास्तविक हल हैं, और हर एक के लिए $1$
लंबाई का एक \omterm{def:g10:numbers:interval}{अंतराल} बताइए जिसमें वह हो। (संकेत: $f(x) = x^3 - 3x + 1$ के परिवर्तन का अध्ययन
कीजिए; उसका \omterm{def:g11:deriv:derivative}{अवकलज} $3x^2 - 3$ है।)
\end{exercise}

\begin{exercise}[$\star\star$]\label{exo:g12:limcont:7}
मान लीजिए $f \colon \intcc{0}{1} \to \intcc{0}{1}$ \omterm{def:g12:limcont:continuity}{संतत} है। दिखाइए कि $f$ का कोई \emph{\omterm{pb:g10:functions:1}{स्थिर बिंदु}} है:
अर्थात् ऐसा $c \in \intcc{0}{1}$ है जिसके लिए $f(c) = c$ हो। (संकेत: $g(x) = f(x) - x$ पर मध्यवर्ती मान
प्रमेय लगाइए।)
\end{exercise}

\begin{exercise}[$\star\star\star$]\label{exo:g12:limcont:8}
एक पदयात्री पहाड़ी रास्ते पर 8:00 बजे चढ़ना शुरू करती है और 20:00 बजे पहुँच
जाती है। अगले दिन वह उसी रास्ते से उतरती है, फिर 8:00 बजे शुरू करके 20:00 बजे
पहुँचती है। दिखाइए कि दिन का कोई ऐसा समय है जब वह दोनों दिन रास्ते के ठीक एक
ही बिंदु पर थी। (संकेत: दोनों स्थिति-फलनों का अंतर लीजिए।)
\end{exercise}

\begin{exercise}[$\star\star\star$]\label{exo:g12:limcont:9}
मान लीजिए $f(x) = x^5 + x^3 + x - 1$ है।
\begin{enumerate}
  \item दिखाइए कि $f$ का अद्वितीय वास्तविक शून्य $\alpha$ है, और $\alpha \in \intoo{0}{1}$ है।
  \item $\alpha$ को $10^{-6}$ तक निकालने की \omterm{thm:g12:limcont:ivt}{द्विभाजन} प्रक्रिया का वर्णन कीजिए, और
        $\intcc{0}{1}$ से शुरू करने पर आवश्यक पुनरावृत्तियों की संख्या बताइए।
\end{enumerate}
\end{exercise}

\section{समस्या: नदी बिना भीगे पार नहीं होती}

\begin{problem}[{सप्ताहांत समस्या --- मध्यवर्ती मान प्रमेय स्थिर बिंदु,
समान-ताप वाले प्रतिध्रुवी बिंदु और स्थिर मेज़ें ज़बरदस्ती पैदा कर देती है:
एक ही विचार से तीन असंभव लगने वाली प्रमेय}]
\label{pb:g12:limcont:1}
किसी देश का नक्शा मसलकर उसी देश की ज़मीन पर कहीं भी गिरा दीजिए: नक्शे का कोई
बिंदु \emph{ठीक} उसी जगह के ऊपर होगा जिसे वह दर्शाता है। हर क्षण भूमध्य रेखा
के दो व्यास-विपरीत बिंदुओं का ताप \emph{ठीक} एक ही होता है। और ऊबड़-खाबड़
ज़मीन पर डगमगाती चौकोर मेज़ को \emph{घुमाकर} हमेशा स्थिर किया जा सकता है। तीन
तमाशे, एक प्रमेय: कोई \omterm{def:g12:limcont:continuity}{संतत} राशि ऋणात्मक से धनात्मक तक बिना लुप्त हुए नहीं जा
सकती (\cref{thm:g12:limcont:ivt})। यह समस्या पहले सीमाओं का अभ्यास कराती है, फिर तीनों चमत्कार
खड़े करती है --- और यह भी ठीक-ठीक जानती है कि यह प्रमेय क्या वादा नहीं करती।

\medskip
\noindent\textbf{भाग I --- सीमाओं में प्रवाह।}
\begin{enumerate}
  \item $\lim_{x \to +\infty} \dfrac{3x^2 - x}{x^2 + 5}$, \ $\lim_{x \to +\infty} \dfrac{2x + 1}{x^2 + 1}$, \ और $\lim_{x \to +\infty} (x^3 - x^2)$ निकालिए (\cref{met:g12:limcont:limits})।
  \item संगत सीमाएँ निकालकर
        $f(x) = \dfrac{3x - 1}{x + 2}$
        की अनन्तस्पर्शियाँ बताइए (\cref{def:g12:limcont:asymptote})।
  \item $\lim_{x \to 3} \dfrac{x^2 - 9}{x - 3}$ (\omterm{def:g10:algebra:expand}{गुणनखंडन} से) और $\lim_{x \to 0} \dfrac{\sqrt{x + 1} - 1}{x}$ (संयुग्मी से) निकालिए।
  \item संपीडन प्रमेय (\cref{thm:g12:limcont:squeeze}) से: $\lim_{x \to 0} x \sin\frac1x$ और $\lim_{x \to +\infty} \frac{\sin x}{x}$ निकालिए।
  \item संयोजन और हावी पदों से: $\lim_{x \to +\infty} \dfrac{\sqrt{4x^2 + x}}{x}$।
\end{enumerate}

\medskip
\noindent\textbf{भाग II --- \omterm{def:g11:quad:discriminant}{मूल}, प्रमाणित।}
\begin{enumerate}[resume]
  \item दिखाइए कि \omterm{def:g10:algebra:equation}{समीकरण} $x^3 + x = 5$ का $\intcc{1}{2}$ में कम से कम एक हल है।
  \item दिखाइए कि $\R$ पर हल अद्वितीय है (\cref{thm:g12:limcont:bijection})।
  \item चिरपरिचित परिणाम सिद्ध कीजिए: \emph{विषम घात के हर बहुपद का कम से
        कम एक वास्तविक \omterm{def:g11:quad:discriminant}{मूल} होता है}। ($p(x) = x^3 + a x^2 + b x + c$ के लिए पूरा तर्क लिखिए: हावी पद
        से दोनों अनंत सीमाएँ निकालिए, फिर किसी बड़े \omterm{def:g10:numbers:interval}{अंतराल} पर मध्यवर्ती मान
        प्रमेय लगाइए।)
  \item प्रश्न 6 के \omterm{def:g11:quad:discriminant}{मूल} को \omterm{thm:g12:limcont:ivt}{द्विभाजन} से ढूँढ़ना: $10^{-6}$ से कम त्रुटि की
        गारंटी के लिए $\intcc{1}{2}$ को कितनी बार आधा करना पड़ेगा (\cref{exo:g12:limcont:9})? \cref{pb:g11:deriv:1} की
        न्यूटन विधि के अंक-दुगुनेपन से तुलना कीजिए: उसे मोटे तौर पर कितने
        क़दम चाहिए होते?
  \item वह प्रतिउदाहरण जो प्रमेय की पहरेदारी करता है: $g(x) = \frac1x$ $g(-1) = -1 < 0$ और $g(1) = 1 > 0$
        संतुष्ट करता है, फिर भी कभी लुप्त नहीं होता। मध्यवर्ती मान प्रमेय
        की कौन-सी परिकल्पना टूट जाती है --- और परिकल्पनाएँ जाँचने के बारे
        में यह उदाहरण क्या सिखाता है?
\end{enumerate}

\medskip
\noindent\textbf{भाग III --- तीन असंभव लगने वाली प्रमेय।}
\begin{enumerate}[resume]
  \item रेखाखंड की स्थिर-बिंदु प्रमेय: यदि $f : \intcc{0}{1} \to \intcc{0}{1}$ \omterm{def:g12:limcont:continuity}{संतत} हो, तो सिद्ध कीजिए कि
        कोई $c$ $f(c) = c$ को संतुष्ट करता है। ($g(x) = f(x) - x$ का दोनों सिरों पर अध्ययन
        कीजिए।)
  \item प्रश्न 11 को मसले हुए नक्शे की भाषा में उतारिए: $f$ की भूमिका कौन
        निभाता है, वह $\intcc{0}{1}$ में (लाक्षणिक रूप से: देश के भीतर) क्यों गिरता
        है, और वह बिंदु कौन-सा है जो ठीक अपनी ही जगह के ऊपर है? (एक विमा
        चुपचाप स्वीकार कर ली गई है --- ईमानदार द्विविमीय कथन ब्राउवर की
        प्रमेय है, जो स्नातक खंडों में है।)
  \item भूमध्य रेखा: मान लीजिए $T$ किसी वृत्त पर \omterm{def:g12:limcont:continuity}{संतत} ताप है, जिसे $T(0) = T(2\pi)$
        वाले $2\pi$-आवर्ती \omterm{def:g11:func:function}{फलन} के रूप में देखा जाए। $g(t) = T(t) - T(t + \pi)$ रखिए। $g(0) + g(\pi)$
        निकालिए, निष्कर्ष निकालिए कि $g$ $\intcc{0}{\pi}$ पर कहीं लुप्त होता है, और
        प्रतिध्रुवी बिंदुओं वाली प्रमेय बताइए।
  \item डगमगाती मेज़: बिलकुल कड़ी चौकोर मेज़ \omterm{def:g12:limcont:continuity}{संतत} (कोई सीढ़ी नहीं!)
        ऊबड़-खाबड़ ज़मीन पर खड़ी है; तीन पाए हमेशा ज़मीन छूते हैं और एक
        $h(\theta)$ ऊँचाई पर लटका रहता है, जो मेज़ के घूर्णन-कोण $\theta$ पर \omterm{def:g12:limcont:continuity}{संतत}
        रूप से निर्भर करती है। समझाइए कि एक चौथाई चक्कर घुमाने पर लटकते
        \omterm{def:g10:numbers:interval}{अंतराल} का चिह्न बदलना क्यों अनिवार्य है --- और निष्कर्ष निकालिए कि
        कोई कोण $\theta^*$ चारों पाए ज़मीन पर टिका देता है। (यह तर्क 2005 में
        एक असली प्रमेय के रूप में छपा था; कैफ़े के मालिक यह तरकीब पहले से
        जानते थे।)
  \item \cref{exo:g12:limcont:8} की पदयात्री ने यही ढाँचा काम में लाया था। प्रश्न 11, 13, 14
        और पदयात्री में साझा ढाँचा तीन पंक्तियों में बताइए --- कौन-सा \omterm{def:g11:func:function}{फलन}
        बनाइए, सिरों पर क्या जाँचिए, और किसका हवाला दीजिए।
  \item मध्यवर्ती मान प्रमेय जो \emph{नहीं} देती: अस्तित्व मुफ़्त में मिला,
        पर वह कौन-से दो प्रश्न खुले छोड़ देती है, और इस समस्या के कौन-से
        औज़ार उनका उत्तर देते हैं (प्रश्न 7 और 9)?
\end{enumerate}

\medskip
\noindent\textbf{भाग IV --- अनन्तस्पर्शियों सहित चित्र।}
\begin{enumerate}[resume]
  \item $f(x) = \dfrac{3x - 1}{x + 2}$ को $3 - \dfrac{7}{x + 2}$ के रूप में फिर से लिखिए और उससे वक्र का पूरा वर्णन
        कीजिए: अनन्तस्पर्शियाँ, हर ओर के परिवर्तन, रेखाचित्र। (एक खिसका हुआ
        \omterm{def:g11:func:reference}{अतिपरवलय} --- \cref{pb:g10:reffunc:1} का औसत-लागत वक्र भी वही था।)
  \item \omterm{def:g11:func:function}{फलन} $f(x) = \dfrac{x^2 - 1}{x - 1}$ $1$ पर परिभाषित नहीं है। $f(1)$ को कौन-सा मान देने से वह
        वहाँ \omterm{def:g12:limcont:continuity}{संतत} हो जाता है --- और विश्लेषक ऐसे बिंदु को \emph{हटाने योग्य}
        असांतत्य क्यों कहते हैं?
  \item महत्तम \omterm{def:g10:numbers:sets}{पूर्णांक} \omterm{def:g11:func:function}{फलन} $x \mapsto \lfloor x \rfloor$ ($\leq x$ वाला सबसे बड़ा \omterm{def:g10:numbers:sets}{पूर्णांक}): वह कहाँ
        \omterm{def:g12:limcont:continuity}{संतत} है, कहाँ उछलता है, और उसके लिए मध्यवर्ती मान प्रमेय का निष्कर्ष
        क्यों टूट जाता है (ऐसा $a < b$ ढूँढ़िए जिसके लिए $\lfloor a \rfloor < \frac12 < \lfloor b \rfloor$ हो पर $\lfloor x \rfloor = \frac12$ का
        कोई हल न हो)?
  \item समापन --- \omterm{def:g12:limcont:continuity}{सांतत्य} के दो चेहरे: स्थानीय वादा (मान $=$ सीमा: कोई
        छलाँग नहीं) और वैश्विक शक्ति (मध्यवर्ती मान प्रमेय: हर मध्यवर्ती मान
        छू लिया जाता है)। भाग III के चारों चमत्कारों को सही चेहरे के नीचे
        रखिए, और नदी पार करने वाले नारे को उसका ठीक-ठीक गणितीय नाम दीजिए।
\end{enumerate}
\end{problem}
